\documentclass[12pt]{article}
\usepackage[utf8]{inputenc}
\usepackage[T1]{fontenc}
\usepackage[margin=1in]{geometry}
\usepackage{setspace}
\usepackage{xcolor}
\usepackage[colorlinks=true, urlcolor=blue!60!black, linkcolor=blue!60!black]{hyperref}
\setlength{\parindent}{2em}
\setlength{\parskip}{1.25em}
\renewcommand{\baselinestretch}{1.0}
\begin{document}
\begin{center}
{\LARGE \textit{Ex parte Milligan}, 71 U.S. (4 Wall.) 2 (1866)\par}
\vspace{0.8em}
{\large Case Brief --- Four-Page Edition\par}
\end{center}
\vspace{1em}

\section{Memory Jogger}

An Indiana civilian sentenced to hang by a Civil War military commission won his release: the Court held unanimously that the commission lacked jurisdiction, and 5--4 that neither the President nor Congress may subject civilians to military trial where the civil courts are open, even while habeas corpus is suspended (\href{https://supreme.justia.com/cases/federal/us/71/2/}{71 U.S. (4 Wall.) 2}).

\section{Detailed Case Facts}

The Habeas Corpus Act of March 3, 1863, authorized President Lincoln to suspend the writ, but required officials to list civilian detainees held in loyal States where federal courts were unimpaired and to discharge any prisoner if a federal grand jury met and adjourned without indicting him. Lincoln suspended the writ on September 15, 1863, for spies, aiders of the enemy, and draft resisters.

Lambdin P. Milligan, a twenty-year Indiana resident who had never served in the military or lived in a rebel State, was arrested at home on October 5, 1864, on General Alvin P. Hovey's order. A military commission in Indianapolis tried him for conspiracy, aiding rebels, inciting insurrection, and violating the laws of war---specifically, helping the secret ``Sons of Liberty'' plan to seize arsenals and free Confederate prisoners. Over his jurisdictional objection, it sentenced him to hang on May 19, 1865, with presidential approval. Meanwhile, the federal Circuit Court sat peacefully in Indianapolis in January 1865, and its grand jury adjourned without indicting him.

\section{Procedural History}

On May 10, 1865, Milligan petitioned the Circuit Court for habeas corpus and discharge under the 1863 Act. The two judges divided on whether the writ should issue, whether Milligan should be discharged, and whether the commission had jurisdiction, and certified the questions to the Supreme Court, which decided the case on April 3, 1866, with opinions delivered in December 1866.

\section{Judicial Votes}

\textbf{9--0} that the commission lacked jurisdiction and Milligan must be discharged; \textbf{5--4} that Congress could not have authorized the commission.
\begin{itemize}
\item \textbf{Majority:} Justice Davis, joined by Justices Nelson, Grier, Clifford, and Field.
\item \textbf{Concurring in the judgment:} Chief Justice Chase, joined by Justices Wayne, Swayne, and Miller.
\item \textbf{Dissents:} None.
\end{itemize}

\section{Holding}

The Indiana commission lacked jurisdiction, and Milligan was entitled to discharge under the 1863 Act. By 5--4, neither the President nor Congress may authorize military commissions to try civilians unconnected to the armed forces in States where the civil courts are open and operating unobstructed, even when habeas corpus is suspended in wartime.

\section{Analysis of Opinions}

\textbf{Davis (majority).} After confirming jurisdiction, Davis rested on four principles. \textit{First}, the Constitution binds rulers and citizens equally in war and peace; no branch may suspend any safeguard except the privilege of habeas corpus. \textit{Second}, every criminal trial exercises the judicial power that Article III vests in life-tenured judges, and the Fifth and Sixth Amendments guarantee indictment and jury trial; the Fifth Amendment's exception for ``cases arising in the land or naval forces, or in the militia, when in actual service'' confirms that civilians keep those rights. \textit{Third}, martial rule cannot rest on threatened invasion or executive discretion; it is permissible only where actual invasion or civil war has closed the courts, and it ends when they reopen. Indiana's federal court was open and trying wartime offenses. \textit{Fourth}, suspending habeas permits detention without judicial release, but ``[t]he Constitution goes no further''; the writ still issues so a court can decide whether the petitioner falls within the suspension, and a lifelong civilian of a loyal State is not a prisoner of war.

\textbf{Chase (with Wayne, Swayne, and Miller).} Milligan had to be released because the 1863 Act required grand-jury proceedings or discharge, impliedly forbidding the President's commission. But the majority went too far in limiting Congress. Under its powers to declare war, raise armies, and govern the forces, Congress may authorize military tribunals for wartime offenses even where courts sit, because courts may be technically open yet ``wholly incompetent'' or sympathetic to rebels. Chase distinguished military law, military government, and martial law proper.

\textbf{Later developments.} \textit{Ex parte Quirin} (1942) upheld commission trials of German saboteurs, including a citizen, as unlawful belligerents, distinguishing Milligan as a non-belligerent. In \textit{Hamdi v. Rumsfeld} (2004), a plurality allowed military detention of a citizen captured abroad, while Justice Scalia's dissent relied on \textit{Milligan}.

\section{Examples: Future Applications}

\begin{enumerate}
\item \textbf{Same-Side (Invalid): Congress Authorizes Commissions for Domestic Saboteurs.} Even with habeas suspended and congressional authorization, trying a non-belligerent civilian by commission while federal courts are open violates \textit{Milligan}'s majority, though Chase's view might save it.
\item \textbf{Same-Side (Invalid): Presidential Martial Law After Riots.} A President acting without statute may not try a student arsonist by military tribunal while courts remain open; even Chase located the power in Congress.
\item \textbf{Opposite-Side (Permitted): Citizen Enlisted With the Enemy.} A citizen who joins a hostile military and sabotages a shipyard out of uniform may be tried by commission under \textit{Quirin}.
\item \textbf{Opposite-Side (Permitted): Citizen Captured Abroad.} A citizen captured bearing arms on a foreign battlefield may be detained under the \textit{Hamdi} plurality, with a meaningful chance to contest his status.
\item \textbf{Fence-Sitter: Open but Compromised Courts.} In an insurrection where threats prevent a federal court from seating juries, Congress authorizes commissions for militia suppliers; Davis's ``unobstructed'' test, Chase's warning, and \textit{Quirin}'s belligerency line leave the result uncertain.
\end{enumerate}

\section{Critique}

\textbf{Post-war courage versus wartime realism.} Critics note that the Court acted only after Appomattox, when ``the public safety [was] assured,'' having declined to intervene during hostilities in \textit{Ex parte Vallandigham}. Chase's concurrence exposes the weakness of a binary ``open courts'' test, since courts may be open yet compromised by intimidation or disloyalty.

\textbf{The \textit{Quirin} tension.} By carving out ``unlawful enemy belligerents'' even for citizens on domestic soil, \textit{Quirin} created a persistent boundary problem between civilians protected by \textit{Milligan} and saboteurs triable by commission.

\section{Key Quotations}

\begin{quote}
``The Constitution of the United States is a law for rulers and people, equally in war and in peace, and covers with the shield of its protection all classes of men, at all times and under all circumstances. No doctrine involving more pernicious consequences was ever invented by the wit of man than that any of its provisions can be suspended during any of the great exigencies of government.''
\end{quote}
\noindent --- Justice Davis, Opinion of the Court, 71 U.S. at 120--21.

\begin{quote}
``Martial rule can never exist where the courts are open and in the proper and unobstructed exercise of their jurisdiction. It is also confined to the locality of actual war.''
\end{quote}
\noindent --- Justice Davis, Opinion of the Court, 71 U.S. at 127.

\begin{quote}
``Not one of these safeguards can the President or Congress or the Judiciary disturb, except the one concerning the writ of habeas corpus. ... The Constitution goes no further. It does not say, after a writ of habeas corpus is denied a citizen, that he shall be tried otherwise than by the course of the common law.''
\end{quote}
\noindent --- Justice Davis, Opinion of the Court, 71 U.S. at 125--26.

\end{document}
